% !TEX program = xelatex
% ============================================================================
% demo 示例报告：单摆法测重力加速度
%   用途：验证"讲义 + data/ → 全套报告 + PDF"的流水线能被跑通。
%   数据为演示数据（见 data/README.md）；正式使用时替换为真实记录与真实信息。
%   版式要点（本工作流 v1.25 的硬规则）：
%     * 每个图/表的声明紧跟在"首次提到它的那句话"之后（同栏内间距 ≤ 60pt）；
%     * 一句话只提一个浮动体；
%     * 正文表用 table[H] + \scriptsize 半栏宽（由 make_data_tables.py 生成，勿手改）。
% ============================================================================
\documentclass{thuemp}
\usepackage{xcolor}
\usepackage{array}
\sisetup{detect-all,per-mode=symbol,separate-uncertainty=true}

% —— 一级标题中文编号（老师小论文格式：一、二、三）——
\ctexset{section={name={,、},number=\chinese{section}}}

\Keyword{单摆；重力加速度；最小二乘拟合；周期测量}

\begin{document}

\emptitle{单摆法测重力加速度}
\empauthor{示例同学}{示例教师}

\fancyhead[CO]{{\footnotesize 示例同学: 单摆法测重力加速度实验报告}}

\twocolumn[
\begin{@twocolumnfalse}
\maketitle

% 数据溯源声明（两处落点：此处 + data/README.md）。
% 注意：不要写成 \newcommand{\databox}[1]{...}——`[1]` 的 `]` 会提前结束
% \twocolumn 的可选参数，报 "Argument of \@newcommand has an extra }"。直接写 \fcolorbox。
\noindent\fcolorbox{black!35}{gray!8}{%
\parbox{\textwidth}{\textbf{数据说明：}本报告为\textbf{流水线演示样例}：\texttt{data/} 中的数据为按 $T=2\pi\sqrt{L/g}$ 反算生成的拟真数据，其中 $g$ 取 $9.79\ \mathrm{m/s^2}$，并叠加了计时扰动，\textbf{不是实测记录}。正式提交前必须以本人实测原始记录替换，并重新完成全部计算、图表与结论。}}
\vspace{4.2em}

\begin{empAbstract}
用单摆法测量重力加速度：改变摆长 $L$ 为 $0.500$–$0.900\ \mathrm{m}$，每次测量 50 个周期的累计时间 $t_{50}$ 以减小计时相对误差。由 $T=t_{50}/50$ 得周期，对 $T^2$–$L$ 作最小二乘直线，斜率 $k=4.03433\ \mathrm{s^2/m}$，相关系数 $R^2=0.999987$，得 $g=4\pi^2/k=9.786\ \mathrm{m/s^2}$；五次独立测量结果的平均值为 $9.788\ \mathrm{m/s^2}$，二者相差 $0.02\%$，与当地重力加速度参考值 $9.79\ \mathrm{m/s^2}$ 的相对偏差为 $0.04\%$。残差无单调趋势，说明 $T^2\propto L$ 的线性模型与数据一致，主要误差来自秒表启停的人为反应时间与摆长计入方式。
\end{empAbstract}

\emptitleEn{Measuring the Gravitational Acceleration with a Simple Pendulum}
\empauthorEn{Demo Student}{Demo Supervisor}
\KeywordEn{simple pendulum; gravitational acceleration; least-squares fitting; period measurement}
\begin{empAbstractEn}
The gravitational acceleration was measured with a simple pendulum. For lengths from $0.500$ to $0.900\ \mathrm{m}$, the accumulated time of 50 periods was recorded to reduce the relative timing error. From a least-squares fit of $T^2$ against $L$, the slope is $k=4.03433\ \mathrm{s^2/m}$ with $R^2=0.999987$, giving $g=4\pi^2/k=9.786\ \mathrm{m/s^2}$, in agreement with the mean of five independent results, $9.788\ \mathrm{m/s^2}$.
\end{empAbstractEn}

\empfirstfoot{2026-01-01}{2026-01-08}{2026000000}{demo@example.com}
\end{@twocolumnfalse}
]

\wuhao
\section{引言}

% 首面含题目/作者/数据说明框/中英文摘要，正文若直接顶到页脚会与 \empfirstfoot 的
% 实验时间/报告时间/学号/E-mail 行重叠（pdf-geo 会报"两行文字重叠"）。
% 模板推荐在第一页正文开头收紧版面高度：
\enlargethispage{-3.3cm}

单摆是测量重力加速度最经典的方法之一\cite{qian2019}。当摆角足够小、空气阻力与悬线质量可忽略时，单摆周期只由摆长与重力加速度决定，因而可通过"测周期、定摆长"反推 $g$。本实验的关键在于两点：一是把计时误差摊薄到多个周期上，二是用最小二乘直线同时检验数据与模型的线性关系，而不是只对单点求平均\cite{lv2015}。
\section{实验原理}

摆长远大于摆球半径、摆角小于 $5^\circ$ 时，单摆的周期为
\begin{equation}
T=2\pi\sqrt{\frac{L}{g}},
\label{eq:period}
\end{equation}
其中 $L$ 为摆长、$g$ 为当地重力加速度。测量 50 个周期的累计时间 $t_{50}$ 后
\begin{equation}
T=\frac{t_{50}}{50},\qquad g=\frac{4\pi^{2}L}{T^{2}}.
\label{eq:g}
\end{equation}
式~\eqref{eq:period} 两边平方得 $T^{2}=\dfrac{4\pi^{2}}{g}L$，即 $T^{2}$ 与 $L$ 成正比。以 $T^{2}$ 对 $L$ 作最小二乘直线 $T^{2}=kL+b$，则
\begin{equation}
g=\frac{4\pi^{2}}{k}.
\label{eq:gfit}
\end{equation}
直线拟合把 5 组数据一次性用上，斜率给出的 $g$ 比单点计算的 $g$ 更稳定，且 $R^2$ 可直接反映数据与模型的线性程度。
\section{实验内容与方法}

固定摆球，其直径 $d$ 由游标卡尺测量一次，改变摆线长度使摆长 $L$ 依次取 $0.500$、$0.600$、$0.700$、$0.800$、$0.900\ \mathrm{m}$。每个摆长下把摆球拉开约 $3^\circ$ 后释放，用光电门或秒表测量 50 个周期的累计时间 $t_{50}$，每个摆长测一次，读数记入原始记录表。

数据处理流程：先由式~\eqref{eq:g} 计算周期与逐次 $g$，再按式~\eqref{eq:gfit} 由 $T^{2}$–$L$ 直线的斜率给出拟合值 $g$；全部过程由脚本 \verb|scripts/generate_plots.py| 完成，原始读数不被修改。
\section{结果与分析讨论}

5 个摆长下的周期、$T^{2}$ 与逐次结果见表~\ref{tab:pendulum-result}。

\input{tables/pendulum-result}

逐次结果在 $9.780$–$9.799\ \mathrm{m/s^2}$ 之间，平均 $9.788\ \mathrm{m/s^2}$，极差 $0.19\ \mathrm{m/s^2}$，说明计时扰动对单次结果的影响约在 $0.1\%$ 量级；把 5 组数据一次性用上、以斜率定 $g$，正是为了压掉这种单次波动。

$T^{2}$ 对 $L$ 的拟合直线见图~\ref{fig:pendulum-fit}。

\begin{figure}[H]
  \centering
  \includegraphics[width=\columnwidth]{figures/pendulum-fit.pdf}
  \caption{单摆 $T^{2}$–$L$ 关系与最小二乘拟合直线，拟合给出 $g=4\pi^{2}/k$}
  \label{fig:pendulum-fit}
\end{figure}

拟合斜率 $k=4.03433\ \mathrm{s^2/m}$、截距 $b=-0.00079\ \mathrm{s^2}$，$R^2=0.999987$，最大残差 $2.52\times10^{-3}\ \mathrm{s^2}$。残差远小于 $T^{2}$ 的量级 $2$–$3.6\ \mathrm{s^2}$，且在 5 个摆长上无明显单调走向，说明 $T^{2}\propto L$ 的线性模型与数据一致，系统误差未表现为趋势项。

由式~\eqref{eq:gfit} 得 $g=9.786\ \mathrm{m/s^2}$，与逐次平均 $9.788\ \mathrm{m/s^2}$ 相差 $0.02\%$，说明两点分别求 $g$ 的结果自洽。与当地重力加速度参考值 $9.79\ \mathrm{m/s^2}$ 比较，相对偏差为 $0.04\%$，落在计时误差允许的范围内。

误差来源：一是秒表启停的人为反应时间约 $0.2\ \mathrm{s}$，折算到单个周期约 $0.004\ \mathrm{s}$，是主要随机误差；二是摆长测量把摆球半径计入方式不一致带来的系统误差，本实验按线长加半径计算，若按线长计会使 $L$ 偏小约 $1\%$、$g$ 相应偏小；三是摆角稍大时式~\eqref{eq:period} 的高阶修正项 $\left(1+\theta_0^{2}/16\right)$ 会使实测周期偏大，$3^\circ$ 时修正量约 $0.02\%$，可忽略。
\section{结论}

以 5 个摆长的 50 周期累计计时数据，用 $T^{2}$–$L$ 最小二乘直线测得重力加速度 $g=9.786\ \mathrm{m/s^2}$，$R^2=0.999987$；逐次测量平均值为 $9.788\ \mathrm{m/s^2}$，与拟合值和当地参考值的偏差均在 $0.05\%$ 以内，表明方法自洽、数据与单摆模型一致。
\renewcommand\refname{\heiti\wuhao\centerline{参考文献}\global\def\refname{参考文献}}
\vskip 12pt
{
\renewcommand{\baselinestretch}{0.9}
\liuhao
\bibliographystyle{gbt7714-numerical}
\bibliography{refs}
}

\section*{附录 A\quad 原始数据}

以下为单摆法测重力加速度的原始记录，即摆长与 50 个周期的累计时间，见表~\ref{tab:pendulum-raw}，未作任何数值修改；其派生结果见表~\ref{tab:pendulum-result}。

\input{tables/pendulum-raw}

\end{document}
